\bookchapter{Глава 16: Система учётных записей и облачные сервисы}

\chapterintro{У моего друга в прошлом году украли телефон. Первой его мыслью было не «сколько стоит телефон», а «смогу ли я восстановить свои фотографии». К сожалению, он не включил облачное резервное копирование. Телефон — вещь заменимая, но данные в учётной записи, если потерять, уже не вернуть. В этой главе мы поговорим о том, как застраховать вашу цифровую жизнь по-настоящему.}

\booksection{I. Сколько всего привязано к вашей учётной записи?}

Многие думают, что учётная запись нужна только для скачивания приложений. На самом деле, за одной учётной записью тянется гораздо больше, чем вы представляете:

\begin{itemize}
  \item \textbf{Контакты}: номера телефонов, которые вы сохраняли годами, а то и десятилетиями.
  \item \textbf{Фото и видео}: от рождения ребёнка до выпускного, от первого свидания до свадебного путешествия.
  \item \textbf{Заметки}: ваши мысли, дневники, пароли от аккаунтов (да, многие хранят пароли именно в заметках).
  \item \textbf{История покупок приложений}: платные приложения и встроенные покупки на сотни юаней.
  \item \textbf{Локатор телефона}: если телефон потерян, это единственный инструмент для определения местоположения, блокировки и удалённого стирания данных.
\end{itemize}

Если вы пользуетесь iPhone, всё перечисленное привязано к вашему Apple ID. Если Android — то к вашей учётной записи Huawei/Xiaomi/OPPO/vivo и учётной записи Google. Эти аккаунты — «прописка» вашей цифровой жизни. Потеряете прописку — и многое станет недоказуемым.

\booksection{II. Три обязательных дела — сделайте прямо сейчас}

\begin{enumerate}
  \item \textbf{Включите двухфакторную аутентификацию}: неважно, iPhone у вас или Android, первым делом зайдите в настройки учётной записи и включите «двухфакторную аутентификацию» или «двухэтапную проверку». Это как добавить засов к вашей двери — одного пароля недостаточно, нужен ещё код подтверждения с вашего телефона. Даже если кто-то угадает ваш пароль в интернете, без вашего телефона он не войдёт.
  \item \textbf{Привяжите надёжный способ восстановления}: двухфакторная аутентификация безопасна, но есть риск — что если вы потеряете свой телефон? Поэтому система предложит вам настроить «ключ восстановления» или «резервную почту/номер телефона». Это единственный способ вернуть доступ к аккаунту в экстренной ситуации. Лучше запишите на бумаге, а не храните в телефоне.
  \item \textbf{Убедитесь, что «Локатор» включён}: эта функция после потери телефона поможет определить местоположение, заблокировать устройство, воспроизвести звук и даже удалённо стереть все данные. Это последняя линия обороны. Многие пропускают этот шаг при покупке телефона, а когда нужно — обнаруживают, что не включили, и уже поздно.
\end{enumerate}

\booksection{III. iCloud: «воздух» для пользователей Apple}

Каждый владелец iPhone рано или поздно сталкивается с экзистенциальным вопросом: что делать, когда место в iCloud заканчивается?

Apple бесплатно даёт 5 ГБ, и, честно говоря, этих 5 ГБ не хватит даже на одну резервную копию. Поэтому большинство людей в итоге получают раздражающее уведомление: «Хранилище iCloud заполнено». У вас есть три варианта:

\begin{itemize}
  \item \textbf{Очистка}: зайдите в Настройки ваше имя iCloud Управление хранилищем и посмотрите, что занимает место. Обычно главные потребители — фото и резервные копии. Можно удалить ненужные старые резервные копии.
  \item \textbf{Расширение}: 50 ГБ за 6 юаней в месяц, 200 ГБ за 21 юань в месяц. Лично я считаю, что 50 ГБ достаточно для большинства — в основном для резервного копирования и хранения миниатюр фотографий.
  \item \textbf{Обходной путь}: если не хотите платить, можно периодически экспортировать фото на компьютер и отключить синхронизацию фото iCloud. Но тогда вы потеряете удобство автоматического резервного копирования фотографий.
\end{itemize}

Важное напоминание: Медиатека iCloud — это не «облачный диск», это «синхронизация». Если вы удалите фото на телефоне, оно удалится и в iCloud. Если вам нужно «сохранить в облаке и очистить телефон», то вам нужна не медиатека iCloud, а резервная копия iCloud или сторонний облачный диск.

\begin{tcbinfo}{iCloud+: три инструмента конфиденциальности, доступные только платным пользователям}
При переходе на любой платный тариф iCloud (50 ГБ / 200 ГБ / 2 ТБ и т.д.) вы автоматически становитесь пользователем\textbf{iCloud+}и получаете три функции защиты конфиденциальности:

\begin{itemize}
  \item \textbf{Private Relay (Частный узел)}: «VPN» от Apple. Трафик Safari проходит через серверы Apple, прежде чем отправиться на сайт, и никто (включая Apple) не может одновременно видеть, кто вы и что вы посещаете. Подходит для защиты конфиденциальности в публичных Wi-Fi. Путь:\texttt{Настройки > ваше имя > iCloud > Частный узел}.
  \item \textbf{Hide My Email (Скрыть e-mail)}: генерирует одноразовые случайные адреса электронной почты при регистрации на сайтах и в приложениях. Письма автоматически пересылаются на ваш реальный адрес. Если больше не хотите получать письма от сайта, просто отключите случайный адрес — настоящий email не раскрывается. Путь:\texttt{Настройки > ваше имя > iCloud > Скрыть e-\allowbreak mail}.
  \item \textbf{Пользовательский домен электронной почты}: можно использовать собственный домен (например, @ваша-фамилия.com) для отправки и получения почты iCloud, подходит для личного бренда или семейного управления почтой.
\end{itemize}

Примечание: Private Relay недоступен в некоторых странах и регионах, пользователям в Китае могут потребоваться альтернативные решения.

\end{tcbinfo}

\begin{tcbwarn}{iCloud Advanced Data Protection (ADP) — настоящее сквозное шифрование}
В конце 2022 года Apple запустила\textbf{iCloud Advanced Data Protection (ADP)}, расширив сквозное шифрование (E2EE) с нескольких категорий данных (здоровье, пароли) на\textbf{большинство данных iCloud}— фото, заметки, резервные копии iCloud, информацию об устройстве и ещё 23 категории данных.

\begin{itemize}
  \item \textbf{Что такое E2EE}: ключи шифрования существуют только на ваших устройствах, их нет даже на серверах Apple. Даже если Apple взломают или правоохранительные органы потребуют данные, Apple не сможет предоставить расшифрованные данные.
  \item \textbf{Путь включения}:\texttt{Настройки > ваше имя > iCloud > Advanced Data Protection}, следуйте инструкциям для настройки контакта восстановления / ключа восстановления (\textbf{обязательно}, потеряете ключ — данные не спасти).
  \item \textbf{Важные замечания}: (1) После включения обязательно надёжно храните ключ восстановления, иначе данные не восстановить при потере аккаунта; (2) Некоторые функции, такие как общие альбомы и общая медиатека iCloud, не участвуют в E2EE; (3) Apple ID, зарегистрированные в материковом Китае, пока не поддерживают ADP, необходимо переключиться на аккаунт другого региона.
\end{itemize}

Краткий совет: если в вашей учётной записи есть важные конфиденциальные данные и регион аккаунта поддерживает ADP, настоятельно рекомендуется включить — это самая надёжная защита облачных данных, доступная обычному человеку.

\end{tcbwarn}

\booksection{IV. Синхронизация и резервное копирование — не путайте}

Эти два слова звучат похоже, но последствия путаницы совершенно разные:

\begin{itemize}
  \item \textbf{Синхронизация}двусторонняя: удалили на телефоне — удалилось и в облаке; удалили в облаке — удалилось и на телефоне. Фото iCloud, контакты, заметки — всё это синхронизация. Плюс — согласованность между устройствами, минус — «удалил — значит, действительно пропало».
  \item \textbf{Резервное копирование}одностороннее: телефон делает «снимок» текущего состояния и сохраняет в облаке. Если потом вы удалили что-то на телефоне, в облачной резервной копии это останется. Резервная копия iCloud, локальная копия на компьютере — это всё резервное копирование. Плюс — данные в большей безопасности, минус — занимает место и не обновляется в реальном времени.
\end{itemize}

\booksection{V. Матрица учётных записей пользователя Android}

У Android ситуация сложнее, чем у Apple, потому что ваши данные могут быть разбросаны по нескольким учётным записям:

\begin{itemize}
  \item \textbf{Аккаунт бренда}(Huawei/Xiaomi/OPPO/vivo): управляет облачным пространством, локатором телефона, магазином приложений, кошельком. У каждого бренда своё облачное решение, обычно базовый тариф бесплатный + платное расширение.
  \item \textbf{Аккаунт Google}: если вы пользуетесь Gmail, Google Фото, Google Контактами — всё это в аккаунте Google. Google бесплатно даёт 15 ГБ, что щедрее Apple, но это пространство общее с Gmail и Drive.
  \item \textbf{WeChat/QQ}: они независимы от системы телефона. Ваша история чатов и избранное в WeChat не хранятся ни в аккаунте бренда, ни в Google — только на серверах WeChat.
\end{itemize}

Поэтому пользователям Android при резервном копировании нужно проверять отдельно: облачную копию бренда, резервную копию Google, резервную копию чатов WeChat — все три обязательны.

\booksection{VI. Семейный доступ: пусть близкие тоже получат выгоду}

Функция Apple «Семейный доступ» позволяет делиться с семьёй покупками приложений, музыкой, фильмами, пространством iCloud, а также местоположением. Можно добавить до 6 человек. Настройка: Настройки ваше имя Семейный доступ. У Huawei и Xiaomi есть похожие функции, которые можно найти в настройках аккаунта — «Семейный доступ» или «Семейное облако».

Очень практичный сценарий: настройте семейный доступ на телефонах родителей, и вы всегда сможете видеть их местоположение (при их согласии). Если пожилой человек потеряется или не будет выходить на связь, у вас будет координаты.

\begin{tcbwarn}{Последствия кражи аккаунта серьёзнее, чем вы думаете}
Если хакер получит ваш пароль, он сможет: удалённо заблокировать телефон и требовать выкуп, просматривать ваши фото и контакты, совершать покупки с вашими платёжными средствами. Хуже того, если к аккаунту привязано слишком много всего, восстановление доступа само по себе станет кошмаром. Поэтому: двухфакторная аутентификация — включите прямо сейчас.

\end{tcbwarn}

\begin{tcbpractice}{Практика: проведите полную проверку вашей учётной записи}
\begin{enumerate}
  \item Откройте настройки телефона, зайдите в раздел учётной записи, убедитесь, что двухфакторная аутентификация включена.
  \item Проверьте, настроены ли «Доверенный номер телефона» и «Ключ восстановления», запишите ключ на бумаге.
  \item Зайдите в iCloud или управление облачным пространством, посмотрите, какие данные синхронизируются, отключите ненужное.
  \item Убедитесь, что «Локатор» включён.
  \item Если у вас Android, проверьте, охвачены ли все три: облачная копия бренда, копия Google, копия чатов WeChat.
\end{enumerate}

\end{tcbpractice}
